% Propietario recuperado íntegro: observador, automedida y Stokes holográfico finito.
% Notación propia del propietario original, preservada localmente para su consumo modular.
\providecommand{\X}{\mathcal X}
\providecommand{\Acal}{\mathfrak a}
\providecommand{\bd}{\partial}
\providecommand{\oneE}{\mathbf 1_E}

\section{Fondements algébriques de la couture de phase et de l'automesure}

\begin{definition}[Géométrie de couture de phase]
Une \emph{géométrie de couture de phase} est un système
\[
\X=(X,\oplus,\otimes,\tau,C_n,L,\mathfrak c,v),
\]
où (X) est un support local fini ; \(\oplus\) et \(\otimes\) sont ses
deux lois internes de transport ; \(\tau:X\to D\) est l'application résiduelle ;
\(C_n\) est un cycle orienté ; \(L\) est un réseau dont l'espace réel
\(W=L\otimes_{\mathbb Z}\mathbb R\) porte un produit scalaire positif ;
\(\mathfrak c\in\mathbb P(D^n)\) est une classe résiduelle distinguée non
triviale ; et \(v\in L\) est son représentant primitif dans la chambre fixée.
Un transport \(\omega\in C^1(C_n;W)\) est \emph{admissible} lorsque le mode
harmonique de sa décomposition appartient à la droite
\(\mathbb Rv\subset W\).
\end{definition}

\begin{theorem}[Décomposition exacte--harmonique sur un cycle]
\label{thm:hodge-cycle}
Pour toute (1)-cochaîne \(\omega\in C^1(C_n;W)\), il existe une décomposition
unique
\[
  \omega=d\Phi+u\oneE,
\]
avec \(\Phi\in C^0(C_n;W)\), normalisée par \(\Phi(0)=0\), et \(u\in W\).
\end{theorem}

\begin{proof}
Si \(e_0,\ldots,e_{n-1}\) sont les arêtes orientées du cycle, on définit
\[
 u=\frac1n\sum_{j=0}^{n-1}\omega(e_j),
 \qquad y_j=\omega(e_j)-u.
\]
Alors \(\sum_jy_j=0\). Les sommes partielles
\(\Phi(k)=\sum_{j=0}^{k-1}y_j\), avec \(\Phi(0)=0\), déterminent une
(0)-cochaîne bien définie et vérifient \(d\Phi(e_j)=y_j\), y compris sur l'
arête finale. Pour l'unicité, de
\(d\Phi+u\oneE=d\Psi+w\oneE\) on obtient, après sommation sur le cycle,
\(n(w-u)=0\) ; donc \(u=w\), et la normalisation élimine la constante
restante de \(\Phi-\Psi\).
\end{proof}

\begin{definition}[Défaut de couture]
Le relèvement continu par morceaux \(\Gamma_\omega:[0,1]\to W\) est fixé par
\[
 \Gamma_\omega\!\left(\frac{j+1}{n}\right)
 -\Gamma_\omega\!\left(\frac jn\right)=\omega(e_j).
\]
Son défaut de couture est
\[
 \bd\Gamma_\omega
 =\Gamma_\omega(1)-\Gamma_\omega(0)
 =\sum_{j=0}^{n-1}\omega(e_j).
\]
\end{definition}

\begin{proposition}[Annulation télescopique du terme exact]
\label{prop:exact-vanish}
Pour toute (0)-cochaîne \(\Phi\), on a
\[
 \sum_{j=0}^{n-1}(d\Phi)(e_j)=0.
\]
Par conséquent, le défaut de couture dépend exclusivement du mode
harmonique \(u\oneE\) de la décomposition du
\cref{thm:hodge-cycle}.
\end{proposition}

\begin{proof}
La somme de \(\Phi(j+1)-\Phi(j)\) est télescopique sur le cycle. Ainsi, si
\(\omega=d\Phi+u\oneE\), alors
\(\bd\Gamma_\omega=nu\).
\end{proof}

\begin{theorem}[Théorème universel de couture]
\label{thm:universal-sewing}
Soit \(\X\) une géométrie de couture de phase et soit \(\omega\) un transport
admissible. Il existe un unique scalaire réel \(\Acal_\X(\omega)\) tel que
\[
 \bd\Gamma_\omega=\Acal_\X(\omega)v,
 \qquad
 \Acal_\X(\omega)
 =\frac{\langle\bd\Gamma_\omega,v\rangle}{\langle v,v\rangle}.
\]
\end{theorem}

\begin{proof}
Par le \cref{thm:hodge-cycle}, \(\omega=d\Phi+u\oneE\), avec \(u\) unique.
L'admissibilité fournit un unique \(\lambda\in\mathbb R\) avec
\(u=\lambda v\). La \cref{prop:exact-vanish} donne
\(\bd\Gamma_\omega=n\lambda v\) ; il suffit de définir
\(\Acal_\X(\omega)=n\lambda\) et de projeter orthogonalement sur
\(\mathbb Rv\).
\end{proof}

\section{État conjoint observateur--observable, automesure intrinsèque et identité holographique de Stokes}\label{sec:observer-holo}

Cette section explicite trois affirmations qui, dans le reste du manuscrit, n'apparaissaient
que sous une forme narrative. La première est que l'\emph{observateur} n'introduit pas l'invariant,
mais sélectionne une polarisation, une trivialisation et une orientation de lecture.
La deuxième est que l'\emph{automesure} n'est pas une métaphore, mais l'unique fonctionnelle linéaire
normalisée compatible avec le mode résiduel de bord. La troisième est que ce que l'on appelle le
\emph{principe holographique} admet, dans ce cadre fini, une formulation exacte par
Stokes discret : le même scalaire peut se lire au bord comme défaut de couture ou dans le
\emph{volume intérieur} comme courbure totale projetée.

D'un point de vue historique, ces trois éléments permettent de réécrire la notion classique de
mesure. En métrologie standard, une mesure est traçable lorsqu'elle peut être reliée à une référence
par une chaîne documentée et ininterrompue de comparaisons ; et, depuis 2019, le SI est défini
par des valeurs fixées de constantes de la nature \cite{NISTTraceability2023,BIPMSI2019}.
Le présent cadre ne nie pas ce schéma : il le déplace vers un niveau antérieur, dans lequel l'
« unité » n'est plus d'abord un étalon extérieur, mais un mode résiduel interne. La mesure cesse alors d'
être une comparaison à un artefact et devient la lecture d'une clôture propre.

\subsection{L'observateur comme polarisation, trivialisation et orientation}

L'interprétation géométrique la plus stable du problème de l'observateur est la suivante. Le support
local possède au moins deux polarisations internes, l'une additive et l'autre multiplicative. Dans le langage
standard de la cinématique quantique finie, les deux polarisations engendrent une extension de
Weyl--Heisenberg dans laquelle la phase centrale enregistre l'ordre des transports
\cite{Tolar2014,Vourdas2006}. L'observateur n'ajoute pas de nouvelle dynamique au système ; il sélectionne un
représentant de la même classe de transport.

\begin{definition}[Observateur de bord]
Soit $\X=(X,\oplus,\otimes,\tau,C_n,L,\mathfrak c,v)$ une géométrie de couture de phase et soit
$\omega\in C^1(C_n;W)$ un transport admissible. Un \emph{observateur de bord} est une paire
\[
\mathcal O=(\varepsilon,\Phi),
\qquad
\varepsilon\in\{+1,-1\},
\qquad
\Phi\in C^0(C_n;W),
\]
qui agit sur $\omega$ par
\[
\omega^{\mathcal O}:=\varepsilon\,\omega+d\Phi.
\]
Le signe $\varepsilon$ encode l'orientation de lecture et la $0$-cochaîne $\Phi$ encode la
trivialisation ou le choix de jauge.
\end{definition}

\begin{remark}
Au niveau local, le choix entre lecteur additif et lecteur multiplicatif doit être compris comme une
polarisation du même support de phase. Au niveau global du bord, ce choix est absorbé dans
la classe de la $1$-cochaîne observée. La partie géométriquement pertinente est donc la classe de
cohomologie et non le représentant particulier.
\end{remark}

\begin{proposition}[Covariance de l'observateur]\label{prop:observer-covariance}
Pour tout transport admissible $\omega$ et tout observateur de bord
$\mathcal O=(\varepsilon,\Phi)$, on a
\[
\Acal_\X(\omega^{\mathcal O})=\varepsilon\,\Acal_\X(\omega).
\]
En particulier, l'automesure est invariante par changement de trivialisation et ne change de signe qu'à l'
inversion de l'orientation.
\end{proposition}

\begin{proof}
En utilisant la définition du défaut de couture et la linéarité, on obtient
\[
\bd\Gamma_{\omega^{\mathcal O}}
=
\sum_{e\in C_n}\bigl(\varepsilon\,\omega(e)+(d\Phi)(e)\bigr)
=
\varepsilon\sum_{e\in C_n}\omega(e)+\sum_{e\in C_n}(d\Phi)(e).
\]
La somme du terme exact est télescopique et vaut $0$ par la
Proposition~\ref{prop:exact-vanish}. Donc
\[
\bd\Gamma_{\omega^{\mathcal O}}=\varepsilon\,\bd\Gamma_\omega.
\]
En projetant sur la droite résiduelle $\R v$ au moyen de la formule du
Théorème~\ref{thm:universal-sewing}, on conclut
\[
\Acal_\X(\omega^{\mathcal O})
=
\frac{\langle \bd\Gamma_{\omega^{\mathcal O}},v\rangle}{\langle v,v\rangle}
=
\varepsilon\,\frac{\langle \bd\Gamma_\omega,v\rangle}{\langle v,v\rangle}
=
\varepsilon\,\Acal_\X(\omega).
\]
\end{proof}

\begin{corollary}[L'observateur ne crée pas l'invariant]\label{cor:observer-does-not-create}
Deux observateurs de bord qui ne diffèrent que par la jauge mesurent le même scalaire de clôture. Par conséquent,
l'observateur fixe une lecture, mais ne fabrique pas le mode résiduel.
\end{corollary}

\begin{remark}
Cette proposition résout le problème de l'observateur au niveau exact auquel le présent cadre
peut le résoudre. L'observateur ne disparaît pas, mais son rôle change : il n'introduit pas l'étalon, mais
sélectionne une polarisation et une orientation pour un invariant déjà engendré par la
structure elle-même. En ce sens, la théorie est \emph{covariante par rapport à l'observateur}.
\end{remark}

\subsection{Unicité de la fonctionnelle intrinsèque d'automesure}

L'étape suivante consiste à formaliser l'idée que le système se mesure lui-même. Il faut pour cela
distinguer la densité locale du mode harmonique et la grandeur globale de clôture du bord.

\begin{definition}[Densité d'automesure et clôture globale]
Soit $\omega$ un transport admissible sur $C_n$. Nous définissons la \emph{densité d'automesure}
par
\[
\mu_\X(\omega):=\frac{1}{n}\,\Acal_\X(\omega)
=\frac{1}{n}\,\frac{\langle \bd\Gamma_\omega,v\rangle}{\langle v,v\rangle}.
\]
De manière équivalente, si
\[
\omega=d\Phi+u\oneE
\qquad\text{avec }u\in\R v,
\]
alors $u=\mu_\X(\omega)\,v$.
Nous appellerons \emph{clôture globale} le scalaire
\[
\Acal_\X(\omega)=n\,\mu_\X(\omega).
\]
\end{definition}

\begin{lemma}[Quotient unidimensionnel des transports admissibles]\label{lem:one-dimensional-quotient}
Soit $\mathcal T_{\mathrm{adm}}\subset C^1(C_n;W)$ le sous-espace des transports admissibles. Alors
le quotient
\[
\mathcal T_{\mathrm{adm}}/dC^0(C_n;W)
\]
est un espace vectoriel réel de dimension $1$, canoniquement engendré par la classe de
$\,v\oneE$.
\end{lemma}

\begin{proof}
Par le Théorème~\ref{thm:hodge-cycle}, tout $\omega\in\mathcal T_{\mathrm{adm}}$ s'écrit de manière
unique sous la forme
\[
\omega=d\Phi+u\oneE.
\]
L'admissibilité impose $u\in\R v$, de sorte que la classe de $\omega$ modulo les termes exacts est
entièrement déterminée par le scalaire réel $\lambda$ tel que $u=\lambda v$. Par conséquent,
le quotient est isomorphe à $\R v$, et donc unidimensionnel.
\end{proof}

\begin{theorem}[Unicité de l'automesure]\label{thm:uniqueness-self-measurement}
Soit $F:\mathcal T_{\mathrm{adm}}\to\R$ une fonctionnelle linéaire telle que :
\begin{enumerate}[label=\textup{(\roman*)}]
\item $F(\omega+d\Phi)=F(\omega)$ pour toute $\Phi\in C^0(C_n;W)$ ;
\item $F(v\oneE)=1$.
\end{enumerate}
Alors
\[
F=\mu_\X.
\]
En particulier, l'automesure est l'\emph{unique} fonctionnelle linéaire normalisée compatible avec le mode
résiduel distingué.
\end{theorem}

\begin{proof}
Par le Lemme~\ref{lem:one-dimensional-quotient}, le quotient
$\mathcal T_{\mathrm{adm}}/dC^0(C_n;W)$ est unidimensionnel et engendré par la classe de
$v\oneE$. Toute application linéaire invariante de jauge se factorise par ce quotient et est
entièrement déterminée par sa valeur sur le générateur. La normalisation
$F(v\oneE)=1$ fixe alors de manière univoque la fonctionnelle. Par définition,
\[
\mu_\X(v\oneE)=\frac1n\frac{\langle nv,v\rangle}{\langle v,v\rangle}=1,
\]
et, de plus, $\mu_\X$ est linéaire et invariante de jauge. Donc $F=\mu_\X$.
\end{proof}

\begin{corollary}[Automesure dans le cas HMT]\label{cor:self-measurement-hmt}
Dans la spécialisation dodécaphasique forte, où $n=12$ et $v=v_\alpha$, on a
\[
\mu_{\mathrm{HMT}}(\omega)=\frac{\alpha}{12},
\qquad
\Acal_{\mathrm{HMT}}(\omega)=\alpha.
\]
Ainsi, le nombre $\alpha$ est la clôture globale du bord et $\alpha/12$ est la densité élémentaire d'
automesure par phase.
\end{corollary}

\begin{remark}
Cette formulation change le sens d'« unité ». En métrologie classique, l'unité est une
référence extérieure à laquelle un résultat est relié par une chaîne d'étalonnages
\cite{NISTTraceability2023,BIPMSI2019}. Ici, la normalisation provient du système lui-même : le
vecteur $v$ n'est pas importé de l'extérieur, mais engendré intérieurement comme relèvement minimal de la classe
résiduelle distinguée. En ce sens précis, l'unité est une \emph{automesure}.
\end{remark}

\subsection{Principe holographique fini : bord, \emph{volume intérieur} et Stokes discret}

Ce qui apparaissait dans le prologue comme une intuition holographique admet désormais une forme exacte. Le fait
crucial est que le même scalaire peut se lire de deux façons : comme holonomie de bord ou comme
courbure totale projetée du \emph{volume intérieur}. Cela n'identifie pas la théorie à l'holographie
quantique-gravitationnelle de 't Hooft ou de Susskind \cite{tHooft1993,Susskind1995} ; cela produit en revanche
un analogue fini et littéral de réduction au bord.

\begin{definition}[\emph{Volume intérieur} fini à bord cyclique]
Un \emph{volume intérieur} fini à bord cyclique est un complexe cellulaire orienté de dimension $2$
\[
K=K^{(0)}\cup K^{(1)}\cup K^{(2)}
\]
dont la frontière orientée satisfait
\[
\partial K\cong C_n.
\]
Étant donnée une $1$-cochaîne $\Omega\in C^1(K;W)$, nous appellerons \emph{courbure discrète} la
$2$-cochaîne
\[
F_\Omega:=d\Omega\in C^2(K;W).
\]
Si $\Omega\vert_{\partial K}=\omega$, nous dirons que $\Omega$ est une \emph{extension au volume intérieur} du
transport de bord $\omega$.
\end{definition}

\begin{theorem}[Stokes discret projeté]\label{thm:discrete-stokes-projected}
Soit $K$ un volume intérieur fini de frontière $\partial K\cong C_n$, et soit $\Omega\in C^1(K;W)$ une
extension au volume intérieur d'un transport admissible $\omega$ sur la frontière. Alors
\[
\bd\Gamma_\omega
=
\sum_{e\subset \partial K}\Omega(e)
=
\sum_{f\in K^{(2)}}F_\Omega(f).
\]
En particulier,
\[
\Acal_\X(\omega)
=
\frac{\left\langle \sum_{f\in K^{(2)}}F_\Omega(f),v\right\rangle}{\langle v,v\rangle}.
\]
\end{theorem}

\begin{proof}
La première égalité est simplement la définition du défaut de couture appliquée à la restriction de
$\Omega$ à la frontière. La seconde est la formule de Stokes discret pour les cochaînes : la somme orientée
sur la frontière d'un complexe coïncide avec la somme, sur toutes les $2$-cellules, du cobord
orienté de la $1$-cochaîne. C'est-à-dire,
\[
\sum_{e\subset \partial K}\Omega(e)=\sum_{f\in K^{(2)}}(d\Omega)(f)=\sum_{f\in K^{(2)}}F_\Omega(f).
\]
En projetant sur la droite résiduelle $\R v$ et en utilisant le
Théorème~\ref{thm:universal-sewing}, on obtient la formule finale.
\end{proof}

\begin{corollary}[Obstruction à un \emph{volume intérieur} plat]\label{cor:no-flat-bulk}
Si un transport de bord admissible $\omega$ admet une extension plate au volume intérieur,
c'est-à-dire une $\Omega$ avec $F_\Omega=0$, alors
\[
\Acal_\X(\omega)=0.
\]
De manière équivalente, si $\Acal_\X(\omega)\neq 0$, toute extension du bord au volume intérieur doit porter
une courbure discrète non nulle.
\end{corollary}

\begin{proof}
C'est une conséquence immédiate du Théorème~\ref{thm:discrete-stokes-projected}. Si
$F_\Omega=0$, alors la somme projetée des courbures vaut $0$, et donc
$\Acal_\X(\omega)=0$.
\end{proof}

\begin{corollary}[Lecture holographique finie]\label{cor:finite-holography}
L'invariant de clôture peut se lire indifféremment comme observable de bord ou comme courbure
intégrée du \emph{volume intérieur} :
\[
\Acal_\X(\omega)
=
\frac{\langle \bd\Gamma_\omega,v\rangle}{\langle v,v\rangle}
=
\frac{\left\langle \sum_{f\in K^{(2)}}F_\Omega(f),v\right\rangle}{\langle v,v\rangle}.
\]
En ce sens précis, le système satisfait un principe holographique fini.
\end{corollary}

\begin{remark}
Le corollaire précédent est la forme la plus littérale et la plus forte du principe holographique que le manuscrit
puisse soutenir par des preuves internes. Il n'affirme pas une dualité gravitationnelle complète. Il affirme quelque chose de plus
modeste et d'exact : le scalaire physico-mathématique pertinent du système \emph{se factorise par le bord} et,
s'il s'étend à un \emph{volume intérieur}, coïncide avec la courbure totale projetée de ce \emph{volume intérieur}.
Lorsque $\Acal_\X(\omega)=\alpha$, l'égalité dit que $\alpha$ est simultanément une holonomie
de bord et une courbure intégrée.
\end{remark}

\begin{corollary}[Conséquence géométrique pour HMT]\label{cor:hmt-bulk-curvature}
Dans le cas HMT fort, tout remplissage du volume intérieur du cycle dodécaphasique qui étend le transport de
bord canonique doit satisfaire
\[
\left\langle \sum_{f\in K^{(2)}}F_\Omega(f),v_\alpha\right\rangle
=
\alpha\,\langle v_\alpha,v_\alpha\rangle
=
54\alpha.
\]
En particulier, le bord HMT n'admet pas de remplissage plat ; son \emph{volume intérieur} porte nécessairement une courbure
totale projetée non nulle.
\end{corollary}

\begin{proof}
Il suffit d'appliquer le Corollaire~\ref{cor:finite-holography} avec $v=v_\alpha$ et d'utiliser le fait que
$\langle v_\alpha,v_\alpha\rangle=54$.
\end{proof}



\begin{remark}[Sur la géométrie du \emph{volume intérieur}]
Le Théorème~\ref{thm:discrete-stokes-projected} n'identifie pas à lui seul la géométrie métrique du
\emph{volume intérieur} ; il impose toutefois une condition négative très forte : lorsque $\Acal_\X(\omega)\neq 0$, le
remplissage ne peut être exactement ou discrètement plat. Si, en outre, on cherche une complétion globale
compatible avec la lecture modulaire du bord, le candidat naturel cesse d'être le disque euclidien et devient
un \emph{volume intérieur} de type modulaire. Dans la théorie classique, les fonctions modulaires vivent dans le
demi-plan supérieur $\mathbb H$ et l'invariant de Klein $J(\tau)$ classifie le tore complexe global
\cite{DLMFModular}. Le présent résultat ne démontre pas encore cette identification, mais justifie
pourquoi l'intuition d'un \emph{volume intérieur} non euclidien est mathématiquement raisonnable.
\end{remark}

\subsection{Conséquences structurelles de la polarisation, de l'automesure et de l'identité discrète de Stokes}

Les trois éléments précédents permettent de fixer, avec un degré de précision auparavant seulement esquissé,
l'interprétation de l'observateur, de l'automesure et de l'holographie.

\begin{enumerate}[label=\textup{(\arabic*)}]
\item L'\textbf{observateur} est un choix de polarisation, d'orientation et de jauge. Il change la lecture,
mais non le mode résiduel, sauf par le signe associé à l'orientation.
\item L'\textbf{automesure} est l'unique fonctionnelle linéaire normalisée sur l'espace admissible des
transports. Elle n'a pas besoin d'étalon extérieur parce que sa normalisation provient du relèvement minimal interne.
\item Le \textbf{principe holographique} acquiert une forme exacte : le scalaire pertinent peut se lire
intégralement au bord ou intégralement comme courbure totale projetée du volume intérieur.
\end{enumerate}

Cette triple identification permet de reformuler plus fortement la thèse centrale du manuscrit :
non seulement $\alpha$ est un invariant archimédien de bord, mais, dans le cadre fini
considéré ici, $\alpha$ est le premier scalaire qu'une cellule auto-inertielle de phase produit en
se mesurant elle-même et en projetant toute sa courbure effective sur le mode résiduel distingué.
